\documentclass[10pt,a4paper]{article}
\usepackage[left=20mm,right=20mm,top=18mm,bottom=19mm,headheight=14pt]{geometry}
\usepackage{fontspec,xeCJK,amsmath,amssymb,booktabs,array,enumitem,xcolor,fancyhdr,titlesec}
\setmainfont{Times New Roman}\setsansfont{Arial}\setCJKmainfont{Songti SC}
\usepackage[hidelinks]{hyperref}
\definecolor{Ink}{HTML}{183B4E}\definecolor{Muted}{HTML}{52616B}\definecolor{Warn}{HTML}{A33A2B}
\setlength{\parindent}{0pt}\setlength{\parskip}{5pt}\setlength{\emergencystretch}{3em}
\setlist[itemize]{leftmargin=1.4em,itemsep=2pt,topsep=2pt}
\titleformat{\section}{\large\sffamily\bfseries\color{Ink}}{}{0pt}{}
\titlespacing*{\section}{0pt}{10pt}{4pt}
\pagestyle{fancy}\fancyhf{}\renewcommand{\headrulewidth}{0pt}\fancyfoot[C]{\small\thepage}
\newcommand{\external}{\textcolor{Warn}{\textsuperscript{*}}}
\begin{document}
\fancyhead[L]{\small\sffamily Abelian sixfolds / Reading guide}
{\LARGE\sffamily\bfseries\color{Ink}A guide to the 21-case proof}\par
{\small English edition / 11 October 2026 / Matched to the Chinese edition}
\section*{The problem and the scope}
For a complex abelian sixfold $A$, the manuscript studies the algebraicity of rational Hodge classes. Its full conclusion is conditional on explicitly stated external assertions. It also proves structural reductions and several case conclusions using the cited classical and geometric results alone. A power $A^n$ has dimension $6n$.
\section*{How the 21 cases arise}
Up to isogeny, first separate simple varieties from products. For a simple sixfold, Albert endomorphism type, center degree and module rank give eleven coarse rows S1--S11. For a nonsimple sixfold, the ten nontrivial partitions of six give N1--N10. These rows are not a list of all Hodge groups: signatures, field embeddings, isogeny repetitions and polarization types can create strata inside a row.
\section*{The shared proof architecture}
\begin{itemize}
\item Identify the \emph{actual joint} Hodge group from its representation, commutant and Hodge cocharacter. Factorwise Hodge conjectures alone do not settle a product.
\item Compute the exterior or tensor invariants. Pairings give divisors; determinant tensors give Weil or mixed exceptional spaces.
\item Algebraize the exceptional generators by a compatible polarization change, a geometric correspondence or an explicit conditional input.
\item Use rational descent and isogeny invariance. Assert all powers only where the required invariant-generation argument is supplied.
\end{itemize}
\section*{What to take away}
The AI edition starts with four pages on polarization and trace constructions, stabilization with algebraic contraction, the one-cycle criterion and tensor contractions, actual Hodge-circle tests, and full-Clifford matrix extraction. They identify reusable mechanisms and their limits; they do not claim that each ingredient is new or that literature priority has been established.
\section*{External assumptions: read before the conditional cases}
\external\ marks a full-row route using one of four unverified inputs: \textbf{W8} split-eightfold Weil algebraicity; \textbf{CM} arbitrary-dimensional CM Hodge algebraicity; \textbf{KS} the specified full even-Clifford Kuga--Satake correspondence; \textbf{Aux} the specified degree-four auxiliary spin tensor. Aux is separate from KS. The official 7 October record withdraws the W8, KS and K3-product manuscripts; their assertions are retained only as hypotheses. CM is not in that three-item list and is not independently verified here. See \href{https://github.com/openai/math/blob/main/history.md}{the official history} and the JSON contracts.
The supplied external review reports ``no major gap'' \emph{under the stated assumptions}; it is not a new verification, a proof of those assumptions, or a formal certification.
\newpage
\section*{Case map: methods and reading effort}
The effort labels are relative reading estimates, not a ranking of mathematical novelty. ``Standard'' still requires the hypotheses of the cited results.
\par\medskip\small\begin{tabular}{@{}p{0.15\linewidth}p{0.67\linewidth}p{0.12\linewidth}@{}}\toprule
Rows & Method & Effort\\\midrule
S1 & Representation and Hodge-weight exclusion; then pairings. & Medium\\[4pt]
S2--S4 & Full symplectic products via commutants and Goursat. & Standard\\[4pt]
S5--S6 & Quaternionic multiplicity spaces; symplectic invariants. & Standard\\[4pt]
S7\external & Weil plane; quaternionic generation; W8 contraction. & High\\[4pt]
S8\external & Unitary dichotomy; balanced Weil plane via W8. & Medium\\[4pt]
S9 & Real-quadratic polarization correction; split Weil theorem. & High\\[4pt]
S10\external & Determinant torus; W8; independent cover examples. & High\\[4pt]
S11 & Rank bound; interpolation trace form; split Weil theorem. & High\\[4pt]
N1 & Prime factor plus elliptic curve; odd-rank adjustment. & Medium\\[4pt]
N2\external & Residual one-cycle criterion; Aux and CM transfer. & Highest\\[4pt]
N3\external & Repeated-factor reduction; residual KS extraction; CM strata. & Highest\\[4pt]
N4 & Actual joint Hodge-circle test; odd-rank adjustment. & High\\[4pt]
N5 & Surface separation; fourfold Weil plane. & Medium\\[4pt]
N6 & Elliptic isogeny classes; fourfold Weil pullbacks. & Medium\\[4pt]
N7--N9 & Direct low-dimensional product theorem. & Standard\\[4pt]
N10 & Elliptic isogeny classes and pairing invariants. & Standard\\[4pt]
\bottomrule\end{tabular}\par\normalsize
\section*{Powers: the precise ledger}
All powers are stated for S1--S11 and N7--N10. For the full N1, N4, N5 and N6 rows this package records only the stated sixfold conclusion. N2 and N3 have full-row sixfold conclusions, with all powers proved for their specified residual models and other explicitly treated strata. This does not yield an unrestricted all-powers theorem for every sixfold.
\section*{Suggested reading paths}
\textbf{Quick orientation:} this guide, then the AI edition's first four pages. \textbf{First complete examples:} S2--S6 and N7--N10. \textbf{Reusable constructions:} S9, S11 and N4; then the one-cycle/all-powers propositions. \textbf{Hard conditional constructions:} the external contracts, then N3 and N2. Use the theorem graph to locate prerequisites and the original edition for wider context. Review reports have separate provenance and do not change a theorem's logical status.
\end{document}
